\chapter{Related Work}
\label{sec:related}

\section{Approaches to text simplification}
\label{sec:related-baselines}

\paragraph{The comparison simplifier} is an off-the-shelf checkpoint\footnote{\hfmodel{eilamc14/bart-large-text-simplification}} already fine-tuned for simplification, against which the model fine-tuned here is evaluated, a core objective of \cref{rq:comparison}.
Its architecture is \gls{bart}-large \parencite[7872]{lewis-etal-2020-bart} at 406\gls{million} parameters \parencite[2]{cohen2026simplify}.
It was fine-tuned not by \gls{bart}'s original authors but by \textcite[1, 4]{cohen2026simplify}, who released it with their comparative study.
Its training data is the pre-cleaned WikiLarge mirror\footnote{\hfdataset{eilamc14/wikilarge-clean}} this project also trained on, with the prefix \texttt{Simplify:} added to every source sentence \parencite[3--4]{cohen2026simplify}.
The mechanism is plain supervised fine-tuning, with no control tokens, no auxiliary objective, and no reranking.

Three limitations affect the evaluation for \cref{rq:comparison}:
Because it shares its training corpus with this project's sentence model, which additionally dropped mojibake rows (\cref{sec:corpus-preparation}), a fractional difference in \gls{sari} between the two says little about either training run (\cref{sec:rq1}).
Because this project passes inputs without that prefix, the comparison simplifier was scored on inputs formatted differently from its training data (\cref{sec:limitations}).
And because it was fine-tuned on isolated sentences, it has no document counterpart, so it appears at one granularity only.

\paragraph{Controllable simplification with control tokens} trains a model to follow an explicit request.
\Gls{access} \parencite[4691]{martin-etal-2020-controllable} is the canonical instance, and its successor \gls{muss} \parencite[1653]{martin-etal-2022-muss} is the baseline \gls{bless} itself uses \parencite[13293]{kew-etal-2023-bless}.
The architecture is a \gls{seqseq} model, in \gls{muss}'s case \gls{bart}-large.
The training setup conditions each training pair on attributes computed from the pair itself, such as compression ratio, paraphrase amount, and lexical and syntactic complexity.
The model thus learns to follow a request instead of producing one fixed style, and at inference the attributes are set to the values the output should have.

\pagebreak

The limitation that matters here is that the control tokens are proxies \parencite[4691]{martin-etal-2020-controllable}: the model is conditioned on a measured attribute of the source--target pair, and nothing at inference checks that the property was achieved in the output.
The prompted \gls{llm} approach of this thesis asks for control through the audience parameter instead of training it in.
It inherits that limitation: nothing in this thesis tests whether the audience parameter helps the audience it names.

\paragraph{Prompted \glspl{llm}} are represented here by \gls{bless} \parencite[13291]{kew-etal-2023-bless}, from which the third approach of this thesis takes its prompt.
It evaluates 44 models, open- and closed-weight, differing in size, architecture, pretraining methods, and accessibility.
The setup is few-shot in-context learning on three test sets from different domains, namely Wikipedia, news, and medical texts, with three stylistically distinct prompts.
Each prompt carries three demonstrations drawn at random from the dataset's validation split \parencite[13293]{kew-etal-2023-bless}.
Outputs are sampled (nucleus sampling with \(\text{\gls{topp}} = 0.9\), temperature~\num{1.0}, at most \num{100}~tokens) under three random seeds \parencite[13293]{kew-etal-2023-bless}.
No model is fine-tuned; the task is given only through the prompt.
\Gls{bless}'s best models perform comparably with dedicated trained simplification baselines \parencite[13291]{kew-etal-2023-bless}.
A prompted model is therefore a demanding comparison for the 139\gls{million}-parameter model fine-tuned here.

Two limitations of \gls{bless} matter here.
It is sentence-only, so its results do not cover documents.
It leaves the study of in-context learning strategies for simplification to future work \parencite[13293]{kew-etal-2023-bless}, which is why the demonstrations here are fixed (\cref{sec:method-three}).

\paragraph{Document-level simplification} supplies both the corpus and the metric on which the document results rest.
\textcite[7999]{sun-etal-2021-document} assembled D-Wikipedia from aligned English and Simple English Wikipedia articles, keeping only the lead section of each.
The authors evaluated it with four baselines (a vanilla Transformer, the context-conditioned sentence simplifier SUC, BertSumextabs, and \gls{bart}) \parencite[8003]{sun-etal-2021-document}.
They also introduced \gls{dsari} to score it \parencite[8002]{sun-etal-2021-document}.

Like the fine-tuned document checkpoint here, SimDoc \parencite[1, 3]{vasquezrodriguez2024simple} fine-tunes a pretrained \gls{seqseq} model on document pairs.
SimDoc uses \gls{tfive}, fine-tuned on professionally created simplification corpora extended into (complex text, simple text, readability label) triples.
A combined loss joins simplification, readability classification, and a coherence reward, steered by task tokens in the input.
The models fine-tuned in this thesis have no coherence objective (\cref{sec:future-work}).

\pagebreak

\section{Browser-based deployment}
\label{sec:related-browser}

The closest deployed system to this one is the on-device Firefox simplifier of \textcite[105]{romero-etal-2025-efficient}.
It runs \gls{bart}-base \parencite[107]{romero-etal-2025-efficient}, quantised for in-browser execution.
It is deployed through \texttt{transformers.js} so that inference runs inside the browser process itself.
The model is trained on a synthetic corpus of \num{2909} examples generated by prompting \glspl{llm}, which its authors report as outperforming the \num{300}\gls{thousand}-example WikiLarge corpus on standard benchmarks \parencite[105]{romero-etal-2025-efficient}.
A single control token sets the target reading level, in place of the multi-token \gls{access} scheme.
However, separate models trained for one reading level each outperform that control-token model \parencite[110]{romero-etal-2025-efficient}.
The authors find that quantisation reduces model size by up to \qty{75}{\percent} with little performance degradation \parencite[105]{romero-etal-2025-efficient}, so a simplification model may be served entirely on a user's device.

Unlike the Firefox simplifier, the extension built here calls a local backend service instead of running the model in the browser process.
Both arrangements keep inference on the user's own machine, which is the axis the privacy argument of \cref{sec:architecture} turns on.
\textcite{romero-etal-2025-efficient} report quality on standard sentence benchmarks, so their work makes no claim about behaviour on authentic pages, which is the question \cref{rq:deployment} asks.

\section{Behavioural robustness beyond held-out accuracy}
\label{sec:related-robustness}

The deployment half of this thesis takes its framing from CheckList \parencite[4902]{ribeiro-etal-2020-beyond}, a behavioural testing methodology.
It borrows from behavioural testing in software engineering, takes as its input a task-agnostic matrix of linguistic capabilities crossed with test types, generates many test cases from templates and lexicons, and reports a failure rate per test.

It distinguishes three test types \parencite[4902--4903]{ribeiro-etal-2020-beyond}: minimum functionality tests for a targeted behaviour, invariance tests for perturbations that should not change the prediction, and directional expectation tests for perturbations whose effect is known in advance.
Held-out accuracy is the standard measure of generalisation, but it often overestimates real performance, because held-out sets often carry the same biases as the training data \parencite[4902]{ribeiro-etal-2020-beyond}.
\textcite[4906, 4910]{ribeiro-etal-2020-beyond} present CheckList as a complement to benchmarks, not a replacement; passing its tests does not imply that a model is free of the tested failure.

The site audit of \cref{sec:deployment-results} is a robustness study in that sense.
The deployment adds a structural constraint to the linguistic ones: a simplification system that edits a live page must write text back into a document whose markup, links, attributes, and interactive elements survive the edit.

\section{Research hypotheses and expectations}
\label{sec:hypotheses}

Four hypotheses follow from the work reviewed above: two on the comparisons that the design of \cref{sec:methodology-i} sets up (\Cref{hyp:finetuning}, \Cref{hyp:llm-vs-finetuned}), one on how the metrics rank the compared outputs (\Cref{hyp:metrics}), and one on the deployment (\Cref{hyp:web}).
All four assume that findings reported for other systems carry over to the ones used here, and each is stated with the result that would refute it.

\begin{hypotheses}
	\item\label{hyp:finetuning} Fine-tuning a pretrained \gls{seqseq} model on task-matched parallel data improves simplification-specific metrics over the same checkpoint before fine-tuning, the one contrast in the grid that holds architecture and scale fixed.
	This is the premise of the fine-tuned systems of \cref{sec:related-baselines}, among them \gls{muss}, which fine-tunes \gls{bart} on mined paraphrases and, in its supervised variant, on labelled simplification data \parencite[1652--1653]{martin-etal-2022-muss}.
	\Cref{hyp:finetuning} fails if the paired bootstrap finds the \gls{sari} or the \gls{dsari} improvement not significant at \(\gls{pval} < 0.05\).
	\item\label{hyp:llm-vs-finetuned} A prompted instruction-tuned \gls{llm} of a recent generation matches or exceeds a fine-tuned model, since \gls{bless} finds that its best models perform comparably with trained simplification baselines \parencite[13291]{kew-etal-2023-bless}.
	At each granularity, \cref{hyp:llm-vs-finetuned} holds if a prompted model scores at or above the fine-tuned checkpoint on every metric compared there (\gls{sari} on ASSET at sentence level; \gls{dsari} and \gls{lens} on D-Wikipedia at document level), fails if every prompted model scores below it on all of them, and remains open if the metrics disagree.
	\item\label{hyp:metrics} Reference-overlap and simplification-specific metrics rank the same outputs differently, because \gls{bleu} leaves latitude for matching the input while \gls{sari} scores additions, kept \glspl{ngram} and deletions against both the source and the references \parencite[404--405]{xu-etal-2016-optimizing}.
	\Cref{hyp:metrics} fails if \gls{bleu} and \gls{sari} order the five sentence-level conditions the same way.
	\item\label{hyp:web} Behaviour on authentic web pages diverges from benchmark behaviour, since held-out accuracy often overestimates real performance \parencite[4902]{ribeiro-etal-2020-beyond} and web pages supply units that no benchmark contains (\cref{sec:background-web}).
	\Cref{hyp:web} is tested qualitatively: it fails if every failure class the site audit finds also appears in the benchmark evaluation.
\end{hypotheses}

\paragraph{How the hypotheses are tested.}
\Cref{hyp:finetuning,hyp:llm-vs-finetuned} are tested across the sentence-level configurations and document-level runs detailed in \cref{sec:methodology-i}.
\Cref{hyp:metrics} is evaluated on the five sentence-level conditions within the same chapter.
Finally, \cref{hyp:web} is tested via the deployment site audit described in \cref{sec:methodology-ii}.
